\documentclass[12pt]{article}
\usepackage[margin=1in]{geometry}
\usepackage{enumitem}
\usepackage{titlesec}
\usepackage{parskip}
\usepackage{graphicx}
\usepackage{xcolor}
\usepackage{hyperref}
\usepackage{fontenc}

\title{\textbf{Mill 1843 RG Notes}\\
\large Mill, John Stuart. \textit{A System of Logic, Ratiocinative and Inductive}.\\
\large Book III, Chapter VIII: Of the Four Methods of Experimental Inquiry
London: John W. Parker, 1843.}
\author{}
\date{}


\begin{document}

\maketitle

\section{Main Idea}



%--------------------------------------------------
\section{General Notes}
%--------------------------------------------------

\begin{itemize}[leftmargin=*, itemsep=4pt]
    \item Five primary approaches to inquiry:
    \begin{itemize}
        \item \underline{Method of Agreement}: If a phenomenon occurs multiple times, and in every case of its occurence there exists only one common factor, then that factor is the cause. (does not guarantee unique causation)
        \item \underline{Method of Difference}: If a phenomenon occurs in one scenario and does not occur in the other, and the two scenarios are identical other than one factor, then that factor is causal. (satisfies necessary \& sufficient, \textcolor{red}{is strongest method available})
        \item \underline{Joint Method of Agreement and Difference}: If multiple instances of a phenomenon have one common factor, while multiple instances where it does not occur are only tied by the absence of that common factor, then that factor is causal
        \item \underline{Method of Residues}: If $A + B + C \rightarrow E_1 + E_2 + E_3$ is observed, and $A \rightarrow E_1$ and $B \rightarrow E_2$ is known, then $C \rightarrow E_3$. (requires prior knowledge of some of the causal chains)
        \item \underline{Method of Concomitant Variations}: If variation in one phenomenon is tied to variation in another phenomenon,then there exists a causal chain between them. (conflates correlation and causation, quite problematic)
    \end{itemize}
    \item Broad focus on eliminating non-causal factors (this also ties to positivity)
\end{itemize}


\section{Important Specific Factual Claims}

\begin{itemize}[leftmargin=*, itemsep=4pt]
    \item
\end{itemize}


\section{Connections}

\begin{itemize}
    \item See Lipset 1959: ``A deviant case, considered within a context which marshals the evidence on all relevant cases, often may actually strengthen the basic hypothesis if an intensive study of it reveals the special conditions which prevented the usual relationship from appearing.'' pg. 70
\end{itemize}


\section{My Critiques}

\begin{itemize}
    \item
\end{itemize}


\end{document}
